\documentclass[a4paper,11pt]{article}
\usepackage[utf8]{inputenc}
\usepackage[french]{babel}
\usepackage[T1]{fontenc}
\usepackage{amsmath,amssymb,amsthm}
\usepackage{geometry}
\geometry{margin=2cm}
\usepackage{enumitem}
\usepackage{hyperref}
\usepackage{booktabs}
\usepackage{array}
\usepackage{xcolor}
\usepackage{listings}
\usepackage{titlesec}
\usepackage{fancyhdr}
\usepackage{lastpage}
\usepackage{graphicx}
\usepackage{etoolbox}
\usepackage{multirow}
\AtBeginDocument{\shorthandoff{;:!?}}

\titleformat{\section}{\large\bfseries\color{blue!60!black}}{}{0pt}{}[\vspace{-0.5ex}\rule{\textwidth}{0.4pt}]
\titleformat{\subsection}{\bfseries\color{blue!40!black}}{}{0pt}{}
\titleformat{\subsubsection}{\itshape\color{blue!30!black}}{}{0pt}{}

\lstdefinestyle{monstyle}{
  frame=single,
  numbers=left,
  numbersep=5pt,
  breaklines=true,
  basicstyle=\small\ttfamily,
  backgroundcolor=\color{gray!5},
  rulecolor=\color{gray!40},
  columns=fullflexible,
  keepspaces=true,
  showstringspaces=false
}
\lstset{style=monstyle}

\pagestyle{fancy}
\fancyhf{}
\fancyhead[L]{\small STA211 -- R\'eseaux de neurones (cours 3)}
\fancyhead[R]{\small Fiche r\'ecapitulative}
\fancyfoot[C]{\thepage/\pageref{LastPage}}
\renewcommand{\headrulewidth}{0.4pt}

\theoremstyle{definition}
\newtheorem*{definition}{D\'efinition}
\theoremstyle{remark}
\newtheorem*{remark}{Remarque}
\newtheorem*{important}{\`A retenir}

\title{\textbf{Fiche r\'ecapitulative -- STA211}\\[4pt]
\large R\'eseaux de neurones -- Cours 3 \\[2pt]
\small Histoire du Deep Learning, architectures modernes CNN
  et perspectives}
\author{Nicolas Thome (CNAM)}
\date{14 juin 2026}

\begin{document}
\maketitle
\thispagestyle{empty}

\begin{abstract}
Ce cours retrace l'histoire du Deep Learning des ann\'ees 80 \`a nos
jours, puis pr\'esente les modules modernes des CNN (ReLU, dropout,
data augmentation, padding), le transfert d'apprentissage (deep
features, fine-tuning), l'adaptation de t\^ache (localisation,
segmentation) et les d\'efis ouverts (apprentissage non supervis\'e,
GANs, VQA).
\end{abstract}

\vfill
\tableofcontents
\newpage

% ===================================================================
\section{Histoire du Deep Learning}
% ===================================================================

\subsection{Ann\'ees 80 : LeNet-5 (d\'ej\`a vu au cours 2)}

LeNet-5 (LeCun et al., 1989) pose les bases des CNN~:
convolution + pooling + FC. $\sim$60\,000 param\`etres, erreur
test 0,95\,\% sur MNIST (0,8\,\% avec distortions).
D\'eploiement pour la lecture de codes postaux.

\subsection{Ann\'ees 90 : hiver du Deep Learning}

\begin{itemize}
  \item R\'eseaux profonds per\c{c}us comme de la \og black
    magic\fg~: bo\^ite noire, manque d'interpr\'etabilit\'e,
    optimisation non convexe difficile.
  \item \^Age d'or des m\'ethodes \`a noyaux~: SVM, kernel trick,
    th\'eorie de la g\'en\'eralisation, optimisation convexe.
\end{itemize}

\subsection{Ann\'ees 2000 : Sac de mots (BoW)}

\begin{itemize}
  \item \textbf{BoW textuel}~: document $=$ histogramme
    d'occurrences de mots, puis classifieur SVM.
  \item \textbf{BoW visuel (BoVW)}~: dictionnaire visuel construit
    par clustering de descripteurs locaux (SIFT) extraits de
    r\'egions d'images.
  \item \'Etat de l'art jusqu'en 2012~: BoW + SVM.
\end{itemize}

\subsection{2006 : renouveau du Deep Learning}

Hinton et al. proposent un apprentissage non supervis\'e pour les
Deep Belief Nets (DBN), utilis\'e comme initialisation avant
r\'etropropagation supervis\'ee. R\'esultats th\'eoriques montrent
l'am\'elioration de la qualit\'e du mod\`ele avec la profondeur.

\subsection{2010 : perc\'ee en reconnaissance vocale}

Premier grand succ\`es du DL sur grands jeux de donn\'ees~:
\emph{Context-Dependent Pre-trained Deep Neural Networks for Large
Vocabulary Speech Recognition} (Dahl et al., 2010).

\subsection{2012 : la r\'evolution AlexNet}

\paragraph{Contexte~:} ImageNet ILSVRC~: 1,2\,M images,
1\,000 classes. \'Evaluation~: top-5 error.

\paragraph{AlexNet (Krizhevsky et al., 2012)}~:
60\,M param\`etres, 650\,K neurones, 630\,M connexions.
\begin{itemize}
  \item 5 couches convolutionnelles + 3 couches FC.
  \item Convolution + ReLU + Pooling.
  \item Apprentissage sur 2 GPU pendant 1 semaine.
  \item Am\'eliorations architecturales~: ReLU (vs sigmo\"ide),
    dropout, data augmentation, overlapping pooling (LRN).
\end{itemize}

\begin{important}
  AlexNet a la m\^eme architecture g\'en\'erale que LeNet, mais
  bien plus grande (60\,M param\`etres vs 60\,K)~: plus de donn\'ees
  et GPU rendent possible l'entra\^inement de grands mod\`eles.
\end{important}

% ===================================================================
\section{Modules modernes des CNN}
% ===================================================================

\subsection{Non-lin\'earit\'es}

\begin{itemize}
  \item \textbf{ReLU (Rectified Linear Unit)}~: $y = \max(0, x)$.
    R\'esout le probl\`eme du gradient vanishing, acc\'el\`ere
    l'apprentissage et la convergence.
  \item \textbf{Leaky ReLU (LReLU)}~: $y = \max(\lambda x, x)$
    avec $\lambda$ pr\'ed\'efini empiriquement.
  \item \textbf{PReLU (Parametric ReLU)}~: $\lambda_k$ appris
    pendant l'entra\^inement.
  \item \textbf{RReLU (Randomized ReLU)}~: $\lambda$ \'echantillonn\'e
    al\'eatoirement pendant l'entra\^inement, fix\'e au test.
  \item \textbf{ELU (Exponential Linear Unit)}~: $y = x$ si $x>0$,
    $y = \lambda(e^x - 1)$ sinon.
\end{itemize}

\subsection{Data augmentation}

Perturbation des images d'entra\^inement~: mirroring, rotation,
stretching, shearing, distortion, perturbation de couleur, etc.
Augmente la taille de l'ensemble d'entra\^inement et la robustesse
aux variations non pertinentes. Appliqu\'ee en train ET en test.

\subsection{Dropout (Hinton et al., 2012)}

\begin{itemize}
  \item Omission al\'eatoire de chaque unit\'e cach\'ee avec
    probabilit\'e $0{,}5$ pendant l'entra\^inement.
  \item Technique de r\'egularisation limitant le sur-apprentissage
    (meilleure g\'en\'eralisation).
  \item Emp\^eche la co-adaptation~: un neurone ne devient utile
    qu'en pr\'esence d'autres neurones sp\'ecifiques.
  \item Interpr\'etation~: moyenne sur de nombreux r\'eseaux.
  \item Au test~: toutes les unit\'es sont utilis\'ees mais les
    poids sont multipli\'es par $0{,}5$ (approximation de la
    moyenne g\'eom\'etrique).
\end{itemize}

\subsection{Padding}

Le \textbf{zero-padding} ajoute des z\'eros autour de l'entr\'ee
pour \'eviter le r\'etr\'ecissement rapide de la dimension
spatiale. Exemple~: padding $p$ sur une image $n\times n$ avec
filtre $f\times f$ \`a stride $1$ donne une sortie
$(n + 2p - f + 1) \times (n + 2p - f + 1)$.

\subsection{Overlapping pooling}

Pooling avec une taille de fen\^etre $z$ sup\'erieure au stride $s$
($z > s$). Utilis\'e dans AlexNet (LRN).

% ===================================================================
\section{Outils et biblioth\`eques}
% ===================================================================

Frameworks disponibles~:
\begin{itemize}
  \item \textbf{Caffe / Decaf}~: script, non modulaire.
  \item \textbf{MatConvNet} (Matlab)~: facile d'utilisation.
  \item \textbf{Torch} (Lua)~: efficacit\'e, modularit\'e.
  \item \textbf{TensorFlow / Theano / PyTorch} (Python)~:
    auto-diff\'erentiation.
  \item \textbf{Keras}~: wrapper haut-niveau sur TensorFlow/Theano.
\end{itemize}

\subsection{Exemple Keras : r\'egression logistique sur MNIST}

\begin{lstlisting}[language=Python]
from keras.models import Sequential
from keras.layers import Dense, Activation
from keras.optimizers import SGD
from keras.datasets import mnist

model = Sequential()
model.add(Dense(10, input_dim=784, name='fc1'))
model.add(Activation('softmax'))

sgd = SGD(lr=0.5)
model.compile(loss='categorical_crossentropy',
              optimizer=sgd, metrics=['accuracy'])

(X_train, y_train), (X_test, y_test) = mnist.load_data()
# pre-processing ...
model.fit(X_train, y_train, batch_size=128,
          epochs=20, verbose=1)
\end{lstlisting}

\subsection{Exemple Keras : CNN}

\begin{lstlisting}[language=Python]
model = Sequential()
model.add(Conv2D(32, kernel_size=5,
          activation='sigmoid', input_shape=(28,28,1),
          padding='same'))
model.add(MaxPooling2D(pool_size=(2, 2)))
model.add(Conv2D(64, (5, 5), activation='sigmoid', padding='same'))
model.add(MaxPooling2D(pool_size=(2, 2)))
model.add(Flatten())
model.add(Dense(100, activation='sigmoid'))
model.add(Dense(nb_classes, activation='softmax'))
\end{lstlisting}

% ===================================================================
\section{Transfert d'apprentissage}
% ===================================================================

\subsection{Probl\'ematique}

Les CNN profonds n\'ecessitent de grands jeux de donn\'ees
annot\'es. Le nombre \'elev\'e de param\`etres rend l'entra\^inement
\emph{from scratch} difficile sur les petits jeux de donn\'ees.

\subsection{Deep Features (transfert pur)}

On extrait les activations d'une couche comme vecteur de
caract\'eristiques~: \textbf{Deep Features} (DF). Ces DF sont
utilis\'ees comme descripteurs \og off-the-shelf\fg{} pour n'importe quelle t\^ache visuelle.

\subsection{Fine-tuning}

On reprend un mod\`ele pr\'e-entra\^in\'e sur ImageNet et on
continue l'entra\^inement (\emph{fine-tuning}) sur le nouveau jeu
de donn\'ees. Meilleure adaptation au domaine cible.

\subsection{R\'esultats exp\'erimentaux}

\begin{center}
\begin{tabular}{lccc}
\toprule
Dataset & Taille & Mod\`ele & Perf. \\
\midrule
\multirow{4}{*}{VOC'07 (20 classes)} & \multirow{4}{*}{$\sim$5\,K ex} & From scratch & 39{,}8 mAP \\
 & & BoW & 61{,}6 mAP \\
 & & Transfert (DF) & 83{,}2 mAP \\
 & & Fine-tuning & 85{,}7 mAP \\
\midrule
\multirow{3}{*}{M2CAI'16 (8 classes)} & \multirow{3}{*}{$6\cdot10^5$ ex} & Transfert & 59{,}3\,\% \\
 & & From scratch & 69{,}1\,\% \\
 & & Fine-tuning & 79{,}1\,\% \\
\bottomrule
\end{tabular}
\end{center}

\begin{itemize}
  \item Petits jeux~: transfert (DF) $>$ Bag of Words.
  \item Jeux moyens~: from scratch possible et comp\'etitif.
  \item Fine-tuning toujours meilleur (petits et grands jeux).
\end{itemize}

\subsection{Adaptation de t\^ache}

\begin{itemize}
  \item \textbf{Localisation (R-CNN)}~: r\'egions propos\'ees,
    extraction de DF, classification. Sur VOC'07~: 58{,}5 mAP
    (vs 34{,}3 pour DPM HoG).
  \item \textbf{Segmentation s\'emantique (FCN)}~: architecture
    encodeur-d\'ecodeur avec convolution 1$\times$1 et
    d\'econvolution (Chen et al., ICLR'15).
\end{itemize}

% ===================================================================
\section{Enjeux actuels et perspectives}
% ===================================================================

\subsection{Apprentissage non supervis\'e}

Le succ\`es du DL est essentiellement supervis\'e. Solution~:
transformer le probl\`eme non supervis\'e en un probl\`eme
supervis\'e~: \textbf{auto-supervision}.

\begin{important}
  \textbf{Generative Adversarial Networks (GAN)} (Goodfellow et al.,
  2014)~: jeu \`a 2 joueurs (g\'en\'erateur vs discriminateur).
  Le g\'en\'erateur apprend la distribution des donn\'ees.
\end{important}

\subsection{Nouvelles t\^aches en IA}

\textbf{VQA (Visual Question Answering)}~: r\'epondre \`a des
questions en langage naturel sur une image. Combinaison vision
+ langage.

\subsection{Limites actuelles}

\begin{itemize}
  \item Compr\'ehension formelle encore limit\'ee.
  \item Optimisation non convexe.
  \item Capacit\'e \`a d\'em\'eler les vari\'et\'es (manifold
    untangling).
  \item Robustesse au sur-apprentissage et g\'en\'eralisation.
\end{itemize}

\subsection{Conclusion}

Le DL a un impact \'enorme sur les r\'esultats exp\'erimentaux,
mais la compr\'ehension th\'eorique et les capacit\'es
d'apprentissage non supervis\'e restent des d\'efis ouverts.

% ===================================================================
\newpage
\section*{Questions cl\'es (QCM)}
% ===================================================================
\addcontentsline{toc}{section}{Questions cl\'es (QCM)}

\begin{enumerate}

  \item Quelle ann\'ee marque la perc\'ee des CNN sur ImageNet~?
    \begin{enumerate}
      \item 2006
      \item 2010
      \item 2012
      \item 2014
    \end{enumerate}
    \textbf{R\'eponse~: c (2012, AlexNet).}

  \item Quel \'etait l'\'etat de l'art en classification image
    avant 2012~?
    \begin{enumerate}
      \item R\'eseaux fully connected
      \item BoW + SVM
      \item R\'eseaux bay\'esiens
      \item For\^ets al\'eatoires
    \end{enumerate}
    \textbf{R\'eponse~: b (BoW + SVM).}

  \item Pourquoi les r\'eseaux profonds ont-ils souffert dans les
    ann\'ees 90~?
    \begin{enumerate}
      \item Manque de donn\'ees
      \item Optimisation non convexe difficile et manque
        d'interpr\'etabilit\'e
      \item Mat\'eriel trop lent
      \item Tous les choix pr\'ec\'edents
    \end{enumerate}
    \textbf{R\'eponse~: d (manque de donn\'ees, optimisation non
    convexe, lenteur, manque d'interpr\'etabilit\'e).}

  \item Quelle fonction d'activation a aid\'e \`a r\'esoudre le
    probl\`eme du gradient vanishing~?
    \begin{enumerate}
      \item Sigmo\"ide
      \item Tanh
      \item ReLU
      \item Softmax
    \end{enumerate}
    \textbf{R\'eponse~: c (ReLU).}

  \item Quel est le r\^ole du dropout~?
    \begin{enumerate}
      \item Acc\'el\'erer l'entra\^inement
      \item R\'egulariser pour \'eviter le sur-apprentissage
      \item Ajouter du bruit aux donn\'ees
      \item Initialiser les poids
    \end{enumerate}
    \textbf{R\'eponse~: b (r\'egularisation par omission al\'eatoire
    de neurones).}

  \item Que sont les Deep Features (DF)~?
    \begin{enumerate}
      \item Les poids de la derni\`ere couche
      \item Les activations d'une couche utilis\'ees comme
        descripteurs
      \item Les gradients de la fonction de perte
      \item Les hyper-param\`etres du r\'eseau
    \end{enumerate}
    \textbf{R\'eponse~: b (activations extraites d'une couche).}

  \item Quelle m\'ethode donne les meilleures performances sur un
    jeu de donn\'ees de taille moyenne ($10^5$ ex)~?
    \begin{enumerate}
      \item Transfert (DF)
      \item From scratch
      \item Fine-tuning
      \item BoW
    \end{enumerate}
    \textbf{R\'eponse~: c (fine-tuning).}

  \item Quel est le principe des GANs~?
    \begin{enumerate}
      \item Un seul r\'eseau qui g\'en\`ere des donn\'ees
      \item Un jeu \`a 2 joueurs~: g\'en\'erateur vs discriminateur
      \item Un r\'eseau qui segmente des images
      \item Un r\'eseau qui traduit des langues
    \end{enumerate}
    \textbf{R\'eponse~: b (jeu \`a 2 joueurs, GAN).}

  \item Quel score R-CNN a-t-il obtenu sur VOC'07 pour la
    localisation~?
    \begin{enumerate}
      \item 34{,}3 mAP
      \item 58{,}5 mAP
      \item 83{,}2 mAP
      \item 95{,}7 mAP
    \end{enumerate}
    \textbf{R\'eponse~: b (58{,}5 mAP).}

  \item Parmi ces \'el\'ements, lequel n'est PAS un module moderne
    des CNN~?
    \begin{enumerate}
      \item ReLU
      \item Dropout
      \item Backpropagation
      \item Data augmentation
    \end{enumerate}
    \textbf{R\'eponse~: c (backpropagation existe depuis les
    ann\'ees 80).}

\end{enumerate}

% ===================================================================
\begin{thebibliography}{9}
% ===================================================================

\bibitem{KSH12}
  A.~Krizhevsky, I.~Sutskever, G.~E.~Hinton.
  \newblock \emph{ImageNet Classification with Deep Convolutional
    Neural Networks}.
  \newblock NIPS, 2012.

\bibitem{HOT06}
  G.~E.~Hinton, S.~Osindero, Y.-W.~Teh.
  \newblock \emph{A fast learning algorithm for deep belief nets}.
  \newblock Neural Computation, 18(7):1527--1554, 2006.

\bibitem{HSK+12}
  G.~E.~Hinton, N.~Srivastava, A.~Krizhevsky, I.~Sutskever,
  R.~Salakhutdinov.
  \newblock \emph{Improving neural networks by preventing
    co-adaptation of feature detectors}.
  \newblock arXiv:1207.0580, 2012.

\bibitem{GPAM+14}
  I.~Goodfellow, J.~Pouget-Abadie, M.~Mirza, B.~Xu, D.~Warde-Farley,
  S.~Ozair, A.~Courville, Y.~Bengio.
  \newblock \emph{Generative Adversarial Nets}.
  \newblock NIPS, 2014.

\bibitem{ARS+16}
  H.~Azizpour, A.~S.~Razavian, J.~Sullivan, A.~Maki, S.~Carlsson.
  \newblock \emph{Factors of transferability for a generic convnet
    representation}.
  \newblock IEEE TPAMI, 38(9):1790--1802, 2016.

\end{thebibliography}

\end{document}
